\documentclass[12pt]{article}
\usepackage[margin=1in]{geometry}
\usepackage{amsmath}
\usepackage{multicol}
\setlength{\parindent}{0pt}

\begin{document}

\begin{center}
{\LARGE Worksheet: Work Done Calculations}\\[0.5em]
{\large Formula: $W = F \times d$}
\end{center}

\vspace{1em}

\textbf{Instructions:} Show all working and don't forget to includes units on your final answers. 

\vspace{1em}

\begin{enumerate}
% --- Basic rearranging practice
\item A force of 12\,N moves an object 4\,m. Calculate the work done.
\item A student pushes a trolley 6\,m using a force of 25\,N. Find the work done.
\item A machine does 1800\,J of work while moving a load 9\,m. Calculate the force.
\item A car engine applies a force of 3000\,N and performs 90{,}000\,J of work. How far did the car move?
\item A force of 40\,N moves an object and does 320\,J of work. How far did it move?
\item A box is pushed 5\,m with a force of 80\,N. Find the work done.
\item A robot arm exerts a force and moves 0.5\,m, doing 150\,J of work. Find the force.
\item A cyclist exerts a force of 160\,N, and the total work done is 960\,J. How far did the cyclist travel?

% --- Unit conversions included
\item A force of 2.5\,kN moves a trailer 6\,m. Calculate the work done.
\item A 500\,mN force moves an object 0.8\,m. Find the work done.
\item A force of 15\,N moves a box 75\,cm. Calculate the work done.
\item A crane performs 15\,kJ of work while lifting a load 3\,m. Find the force exerted.
\item A person lifts a 0.25\,kN object through 1.8\,m. Calculate the work done.
\item A motor performs 4800\,J of work using a 600\,N force. How far did it move the object?
\item A force of 1.2\,kN moves a cart 0.5\,m. Find the work done.
\item A small motor does 60\,J of work while moving a toy car 40\,cm. Find the force exerted.

% --- Trickier composite or contextual ones
\item Two people push a crate together with forces of 120\,N and 150\,N over 3\,m. Calculate the total work done.
\item A 200\,N force pulls a sled 8\,m uphill while friction resists with 50\,N. Find the net work done.
\item A machine applies three equal forces of 400\,N each to move an object 0.6\,m. Determine the total work done.
\item A student lifts a 50\,N weight vertically up 2.0\,m, then lowers it 1.5\,m. Calculate the total work done \emph{by the student}.
\end{enumerate}
\newpage
\begin{center}
{\LARGE Answer Key: Work Done Calculations}
\end{center}
\vspace{0.5em}

\begin{enumerate}
\item $W = 12\times 4 = \mathbf{48~J}$
\item $W = 25\times 6 = \mathbf{150~J}$
\item $F = \dfrac{1800}{9} = \mathbf{200~N}$
\item $d = \dfrac{90000}{3000} = \mathbf{30~m}$
\item $d = \dfrac{320}{40} = \mathbf{8.0~m}$
\item $W = 80\times 5 = \mathbf{400~J}$
\item $F = \dfrac{150}{0.5} = \mathbf{300~N}$
\item $d = \dfrac{960}{160} = \mathbf{6.0~m}$

\item $W = (2.5~\text{kN})(6~\text{m}) = 2500\times 6 = \mathbf{15000~J}~(15~\text{kJ})$
\item $W = (500~\text{mN})(0.8~\text{m}) = 0.5\times 0.8 = \mathbf{0.40~J}$
\item $W = 15\times 0.75 = \mathbf{11.25~J}$
\item $F = \dfrac{15~\text{kJ}}{3~\text{m}} = \dfrac{15000}{3} = \mathbf{5000~N}~(5~\text{kN})$
\item $W = (0.25~\text{kN})(1.8~\text{m}) = 250\times 1.8 = \mathbf{450~J}$
\item $d = \dfrac{4800}{600} = \mathbf{8.0~m}$
\item $W = (1.2~\text{kN})(0.5~\text{m}) = 1200\times 0.5 = \mathbf{600~J}$
\item $F = \dfrac{60}{0.40} = \mathbf{150~N}$

\item $F_{\text{total}} = 120+150=270~\text{N},\; W = 270\times 3 = \mathbf{810~J}$
\item $F_{\text{net}} = 200-50=150~\text{N},\; W_{\text{net}} = 150\times 8 = \mathbf{1200~J}$
\item $F_{\text{total}} = 3\times 400=1200~\text{N},\; W = 1200\times 0.6 = \mathbf{720~J}$
\item Lift: $W_1=50\times 2.0=\;100~\text{J}$; Lower: $W_2=-50\times 1.5=-75~\text{J}$;\\
Total work by student: $\mathbf{25~J}$.
\end{enumerate}
\end{document}
